\chapter{The USSR and the World (1927--1929)}

For nearly two years after the break with Britain in May 1927, and the collapse of the Chinese revolutionary movement and of the Soviet involvement in China, Soviet foreign relations were in the doldrums. Successive approaches from Moscow were ignominiously rebuffed by the British Government. Negotiations with France on debts and credits broke down; and the French Government, though it did not sever diplomatic relations, found a pretext for demanding the recall of Rakovsky, the Soviet Ambassador. Relations with Germany had been temporarily disturbed by her signature of the Locarno treaty and entry into the League of Nations; angry interludes from time to time marked their uneven flow. But, resting on the firm foundation of the secret military agreements, of Germany's desire to avoid an exclusively western orientation, and of common hostility to Poland, they remained closer and more fruitful than Soviet relations with any other country. Relations with Poland had further deteriorated since Pilsudski's coup of May 1926, being embittered by fear that Pilsudski would serve as a willing instrument of the anti-Soviet designs of the western Powers, and later by a persistent, though unsuccessful, Polish attempt to organize the other western neighbours of the USSR---Finland, the Baltic States and Rumania---in a pact or alliance under Polish leadership. Relations with Japan were rendered uneasy by ambiguous Japanese policies in China, and by Japan's jealous control over Manchuria, exercised through her protege, the Chinese warlord Chang Tso-lin. 
% 注：修复了断字 "Nego- tiations" 和 "unsuccess- ful"，修正 "into, the" 为 "into the"，并修正了 OCR 拼写错误 "Chang Tso-Hn" 为 "Chang Tso-lin"。

Paradoxically, the one important initiative undertaken at this time by Soviet diplomacy was participation in international activities at Geneva. Hitherto Soviet cooperation with the League of Nations had been confined to a tenuous link with its health organization. The League had always been denounced in Moscow as an integral part of the oppressive Versailles peace settlement of 1919, and a hypocritical cloak for Allied military preparations. This ban still remained. But, now that Germany had entered the League, absence from the scene of action increased the sense of isolation. In May 1927 a large Soviet delegation arrived for the first time in Geneva to attend the World Economic Conference. Soviet delegates made their mark both in plenary sessions and in the commissions set up by the conference, castigating capitalist procedures and defending the monopoly of foreign trade, but calling for ``peaceful co-existence of the two economic systems''. The absence of concrete results did not remove the impression on both sides that contacts had been established which were susceptible of further development. 

More sensational was the appearance in Geneva, six months later, of a Soviet delegation headed by Litvinov, the deputy Commissar for Foreign Affairs, at the session of the Preparatory Commission for Disarmament.  Litvinov stole the limelight by putting forward a proposal for the total abolition of all military, naval and air armaments. It was a sensational and embarrassing gesture. The commission hastily adjourned the discussion till its next session in March 1928. On this occasion the indefatigable Litvinov submitted a revised plan for total disarmament by stages. When this, too, was shelved, he substituted an alternative draft for a limitation of armaments, which, though less Utopian than its two predecessors, went far beyond anything contemplated by the western Powers. It was greeted with some sympathy only by Germany, whose armaments had been rigidly restricted by the Versailles treaty, and Turkey, a new member of the commission. These proceedings continued to embarrass a majority of the delegates, who found no other resource but a further adjournment, and won for Litvinov and the USSR much favourable publicity in radical circles in the western countries interested in disarmament, which were already impatient at the slow progress made by the commission. 

A further significant step in Soviet relations with the west was taken in the summer of 1928. Kellogg, the American Secretary of State, proposed the signature of an international pact to renounce war ``as an instrument of national policy''. The USSR was not among the fifteen countries originally invited to participate. But, when on the day of the signature, August 27, 1928, an invitation was sent to the USSR, as well as to all other non-signatory countries, to accede to the pact, it was promptly and warmly accepted.  Moreover, when ratification of the pact by the western Powers was delayed, the Soviet Government proposed to its immediate neighbours to conclude a pact to bring the provisions of the Kellogg pact immediately into force between themselves. This subsidiary pact was signed in Moscow on February 9, 1929, amid a blaze of publicity, by the USSR, Poland, Latvia, Estonia and Rumania; Lithuania, Turkey and Persia acceded later. 

In these manoeuvres the hand of Litvinov was clearly visible. Litvinov had now virtually supplanted Chicherin as People's Commissar for Foreign Affairs, though he did not formally succeed to the title till 1930. Chicherin, the eccentric scion of an ancient family who had joined the party, won Lenin's confidence. But strong mutual antipathies divided him from Stalin, who preferred Litvinov's cruder, brusquer style. In 1928, Chicherin, now a sick man, withdrew from active life. The significance of the change was that while Chicherin mistrusted the western countries, especially Britain, from which he had been ignominiously deported in 1918, and felt at home only in Germany, Litvinov had spent many years in Britain, spoke fluent English and had an English wife. For some years Litvinov worked not unsuccessfully, within the limits of Soviet policy, for a rapprochement between the USSR and the western world. 

Ever since the general strike of 1926, the speeches of prominent British politicians had helped to build up the image of Britain prevalent in Moscow as the most implacable enemy of the USSR. The attitude of the Conservative government during this period was inspired by deep mistrust of the USSR and desire to have as few dealings as possible with Moscow. But by the end of 1928 the policy of the cold shoulder had yielded no results, and a milder climate set in. At a time when the Americans and the Germans were beginning to outstrip the British in modern industrial technology, the loss of British markets in the USSR was alarming; and this was attributed to the rupture of relations between the two countries. At the end of March 1929 a party of British businessmen, eighty strong, set out on a visit to the USSR, was enthusiastically received, and collected some orders. A general election in Britain was now imminent. Both the Labour and Liberal Parties inscribed a resumption of relations with the USSR in their platforms. The Labour Party emerged as the strongest of the three parties, and formed a government which carried out its pledge. Full relations were once more established, after some delay, towards the end of the year. But the reconciliation was no more than skin-deep, and did not break the underlying tension between the USSR and the western world. 
% 注：本段清除了断裂的 "I break off" 等杂点。

Relations with the United States were ambivalent, and had a special character of their own. The Soviet leaders recognized that Britain was being rapidly eclipsed by the United States as the major capitalist Power; some expected this to result in acute animosity between the two English-speaking countries. But, in spite of the uniform hostility towards the USSR displayed by the American Government, by the American Federation of Labour, and by the American press, the reaction in Moscow was surprisingly mild, and was tinged with admiration and envy of the accomplishments of American industry. The United States was the most advanced country in the world in industrial technology, in mass production and in standardization; organization in large units of production brought it nearer than any other country to Soviet conditions and requirements. Reliance on American machinery and equipment was an important factor in Soviet industrialization policy; from 1927 onwards the United States began to challenge Germany as the principal supplier of industrial products to the USSR. 
% 注：去除了 "industrializa- ition", "[to challenge" 的乱码。

Still more significant was the extensive employment of American engineers. The staffing of Soviet factories and mines with qualified engineers and technicians had from the first presented a problem. Many who had worked in this capacity before the revolution had disappeared: the loyalty of others was suspect. There were few facilities for training a new generation to take their place. Many German engineers worked in Soviet industry in the early years. But, with the introduction of the first five-year plan and the rapid expansion of Soviet industry, the demand for expertise at the highest levels grew; and it was for the most part American engineers who filled the gap. Dneprostroi was only the first of several major construction projects planned and directed by American chief engineers who brought their staffs with them. The Soviet authorities protected them against outbursts of jealousy on the part of their Russian colleagues, and put more trust both in the efficiency and in the loyalty of American, than of old-style Russian, engineers. By 1929 several hundred highly qualified American engineers were working in the USSR. The number was said to be ``utterly insufficient'', and was soon to be further increased. While the hostility of official circles in the United States to Moscow remained obdurate, a break-through had been effected on the industrial and commercial front. 
% 注：去除了 "jiiindustry", "|/|and", "lithe gap.", "11 circles", "// a break-through", ">11 commercial" 等极为密集的左侧边栏噪点。

The activities of Comintern, being governed in all major questions by the same party leaders who directed Soviet policy as a whole, reflected the uneasiness and ambiguity of Soviet foreign relations in this period. In 1927 the ``united front'', meaning the cooperation of communists with other Left parties or groups in capitalist countries for defined objectives, still figured prominently in the directives of Comintern. But in that year the two most widely publicized experiments in united front tactics, the alliance of the Chinese Communist Party with Kuomintang and the Anglo-Russian trade union committee, both came to grief in conditions of some ignominy (see pp. 93, 102--103 above). The partners with whom cooperation had been sought in these enterprises were denounced as traitors; and the united front in the sense hitherto given to the term was tacitly abandoned. The rift occurred at the moment of a sharp turn for the worse in Soviet relations with the western Powers, when fear of war obsessed the Soviet leaders; and a turn to the Left in Comintern seemed a natural result of the breakdown of conciliatory tactics in Soviet diplomacy, as well as in the relations of communist parties to other Left parties in capitalist countries. That Stalin, having defeated the united opposition, was now moving to the Left in domestic policy, and preparing to confront Bukharin and the Right deviation, was a coincidence which fitted into the same picture. 

From 1928 onwards the new line dominated the proceedings of Comintern. Recognition that the capitalist countries had achieved a phase of ``stabilization'', however ``temporary'', ``relative'' and ``unstable'', was rarer and more grudging. Class antagonisms were becoming more acute; ``class against class'' was a slogan of the new period. The united front was interpreted as a ``united front from below'', meaning cooperation with the rank and file of socialist and social-democratic parties to overthrow their corrupt and traitorous leaders. The sixth congress of Comintern meeting in July 1928---its first for four years, and its longest ever---recorded three periods in its history. The first covered the acute revolutionary ferment of 1917-1921, the second the recovery of capitalism from 1921 to 1927. The third period inaugurated by the congress was one in which the ever-growing contradictions of capitalism heralded its imminent decay, and opened fresh prospects of revolution. The worst enemies of communism were now the temporizing social-democrats. The German delegate bluntly called them ``social-Fascists''. The resolution of the congress admitted that they had some points of contact with the ideology of Fascism; and the new programme of Comintern adopted by the congress bracketed social-democracy and Fascism as twin agents of the bourgeoisie. While the congress was in session, Litvinov was cautiously steering the Soviet Government towards accession to the Kellogg pact, which was announced before the congress ended. No delegate of the Russian party at the congress mentioned the pact. But it was attacked by several delegates of other parties, and in the communist press of western countries, as a hypocritical mask for imperialist aggression; and a resolution of the congress, without mentioning the pact, contained an ironical reference to ``abolition of war'' (in inverted commas) as an example of the ``official pacifism by which capitalist governments mask their manoeuvres''. Any apparent discrepancy between governmental and Comintern policy was probably explained by uncertainties, or unresolved differences of opinion, among Soviet leaders. But in the event the two lines were pursued side by side, and Narkomindel and Comintern went their respective ways without any sense of incompatibility between them. 

Bukharin's downfall was an incidental factor in the proclamation of the ``third period'' in 1928. His quarrel with Stalin turned primarily on economic affairs. But his tenure of office in Comintern had been associated with the conciliatory policies of the united front; and after his disgrace the line swung all the more violently in the opposite direction. An ``objectively revolutionary situation'' was diagnosed in the principal capitalist countries, even before the onset of the world economic crisis lent some slight plausibility to the claim. Revolutionary class war was the primary duty of all communist parties. The term ``social-Fascist'', invented in Germany, was now applied to all ``reformist'' parties of the Left; to seek or tolerate any compromise with them was to be guilty of ``opportunism'' and a ``Right deviation''. These injunctions embarrassed the communist parties of western Europe. In Britain and France they did not prevent some covert communist support for Labour and Socialist candidates in elections. It was in Germany that they were applied with the greatest rigour and with the most disastrous results. The support given by German social-democrats to the Locarno treaty and to the western orientation in German policy earned them the implacable hostility of the Soviet Government and of Comintern. The rift between the German Communist and Social-Democratic Parties persisted, and later proved too deep to be healed even by the imminent danger of Hitler's take-over. 

The breach with other Left parties dealt a fatal blow at the practice of organizing international ``fronts'', in which non-communists of the Left were invited to cooperate with communists in support of causes of common concern (see p. 90 above). Münzenberg, the assiduous and versatile German communist who promoted and directed these joint undertakings, found it necessary, at the sixth congress of Comintern, to plead that such activities had ``nothing in common with an `opportunist policy' or a `Right' deviation''. But they were difficult to reconcile with the uninhibited abuse of social-democrats which was now the rule in communist parties. Nothing short of rigid adherence to the Comintern line was acceptable. Even the once spectacularly successful League against Imperialism wilted in the new climate; and it proved impossible to revive the spontaneous enthusiasm generated at its founding congress in 1927. When a second---and last---congress of the League met two years later in Frankfurt, it was entirely dominated by a Soviet delegation, and the non-communist sympathizers failed to appear, or drifted away. The international society of Friends of the Soviet Union, founded at the tenth anniversary celebrations of the revolution in Moscow in November 1927, was equally short-lived, surviving longer in Britain than elsewhere. The last of these ostensibly non-party demonstrations under Comintern auspices was an Anti-Fascist Congress held in Berlin in March 1929. 

A corollary of the new hard line in Comintern was the enforcement of stricter discipline on communist parties. Ever since 1924, when the Bolshevization of foreign parties was proclaimed as a goal, Comintern had sought from time to time to influence the selection of leaders of these parties. After 1928 such intervention became direct and constant. In the autumn of that year, the central committee of the German party, following a financial scandal, decided to remove its leader, Thälmann, who had owed his rise largely to support from Moscow. The Comintern authorities vetoed the decision and secured its reversal. Early in 1929, Comintern settled a long-standing split in the Polish party by installing in the leadership the group most obedient to its behests; and the existing leaders of the American party were abruptly expelled after a personal intervention by Stalin. Similar changes were rather more cautiously effected in the French and British parties. A feature of most of these changes was a new emphasis in the choice of leaders of unimpeachably working-class origin---Thälmann in Germany, Thorez in France, Pollitt in Britain---which seemed to accord with the Left inclination now prevailing in Comintern, and was a reaction against the trouble caused in the past by dissident intellectuals. Workers generally proved in this respect more malleable. While the new leaders were regularly hailed as Leftists, and those whom they replaced denounced as Rightists, the essential touchstone of the new appointments was prompt and unfailing submission to directives from Moscow. 
% 注：修正 "Thalmann" 的拼写为 "Thälmann"，去除了 "[they replaced" 的括号。

This, however, created another dilemma. The decisions of Comintern were effectively the decisions of the Russian party. They could be, and were, imposed on the foreign parties, but at the cost of alienating a larger and larger number of workers in the countries concerned, who failed to respond to what seemed the wilful and sometimes plainly inappropriate dictates of a remote and alien power. At the end of the nineteen-twenties the communist movement in western countries was declining in numbers and influence, and attracting fewer sympathizers. The British and American parties had no mass following. In Germany, France and Czechoslovakia mass communist parties had split the workers' movement, but never succeeded in dominating it. Everywhere the strengthening of the links that bound the party leaders to Moscow weakened their hold on the rank and file of the workers. It was only when policies in Moscow were radically altered in the middle nineteen-thirties that these losses were retrieved. 
% 注：去除了 "ui^Ui", "dictSLtes" (修正为 dictates), "Ij tries", "I no mass", "IvilU" 等极其破碎的边缘识别字符。

The most important event in Soviet foreign relations in the latter half of 1929 occurred in the Far East. For two years after the debacle of 1927 the Soviet Government was excluded from any participation in Chinese affairs; and the Chinese Communist Party was reduced to scattered underground groups in a few principal cities. In December 1927 the rump of the party, prompted from Moscow, desperately attempted a military coup in Canton. It was a dismal failure, and led to a further massacre of communists and their supporters. About this time the communist peasant leader, Mao Tse-tung, and a communist general, Chu Teh, collected a small force of fugitives and landless peasants, a few thousand strong, in a remote and inaccessible mountain area of south-western China;  and a year later they began to establish their authority over the neighbouring countryside, setting up peasant Soviets. Mao professed formal loyalty to the party and to Comintern. But he went his own way, and had few communications with the party leaders, who mistrusted a movement which pinned its hopes of revolution on peasants in the country and not on workers in the town. Meanwhile Chiang Kai-shek, who had abated nothing of his hostility to Chinese communists and to the USSR, extended the authority of the nationalist government in Nanking over the greater part of China. Chang Tso-lin, the warlord of Manchuria, was killed in the summer of 1928; and at the end of that year Chiang made an agreement with Chang's son and successor for the reunification of China under the Kuomintang flag, while preserving the autonomy of the northern provinces. 
% 注：修正 "debacle of 19275" 为 "1927"，去除了 "succes- sor" 的跨行断字。

The northern provinces abutting on Soviet territory had long been a source of anxiety in Moscow. Here the Chinese Eastern Railway (CER), built and owned by the pre-revolutionary Russian Government on Chinese territory, had long constituted a bone of contention between the two countries (see p. 99 above).  Diplomatic agreements for Chinese representation on its board did not prevent a succession of crises over the control of the railway. But relative calm had prevailed for three years when, in the spring of 1929, the Chinese authorities launched a number of minor attacks on the line. On May 27 they raided the Soviet consulate in Harbin, the headquarters of the CER, arresting officials and seizing papers---a small-scale replica of the raid on the Soviet Embassy in Peking two years earlier. Pronouncements from Nanking left little doubt that the attacks had been inspired by Chiang Kai-shek, and were a first step towards the taking over of the railway. Finally on July 10 the Chinese authorities seized railway installations, closed the Soviet trade delegation and other Soviet establishments in Manchuria, arrested the Soviet general manager of the railway, and expelled him, together with 60 other Soviet officials, from Chinese territory. The Soviet Government, having protested in vain against these high-handed measures, withdrew its staff from the CER, suspended railway communications with China, and demanded the recall of all Chinese officials from the USSR. 

Chiang Kai-shek had assumed that, as happened in 1927, the Soviet Government would protest loudly, but would and could do nothing. This was a grave miscalculation. Soviet interests in central China had never been substantial, and nothing could be done to defend them; the defeat of 1927 was humiliating, but not materially disastrous. But to forfeit Russia's historical position in Manchuria, and to abandon the railway which had been built with Russian capital by Russian engineers, and which was the direct link with Vladivostok, the one Soviet Pacific port, would have been a staggering blow. Moreover, the Red Army had now been built up into an effective fighting force. It was not equipped for a major war. But, once Japan had signified her neutrality in the present quarrel, it was more than a match for the ill-found, ill-disciplined Chinese levies which had fought one another on Chinese soil. Chiang appears also to have assumed that the western Powers would be as favourable in 1929 to his moves against the USSR as they had been two years earlier. This was another miscalculation. Fears of communism had abated, and the British Labour government was just about to resume relations with the USSR. Chiang's aggressiveness looked like another case of something with which the western Powers were all too familiar---the violation by Chinese warlords of foreign treaty rights; and for the first time western sympathies veered to the Soviet side. 
% 注：去除了 "(/the", "I could", "jRussia's", "I ill-disciplined", "-f pathies" 等侧边乱码。

The Soviet Government steadily refused to negotiate on any basis other than the total withdrawal of the measures taken on July 10 and the restoration of Soviet rights over the CER. In August, Blyukher was appointed commander of a reinforced eastern army. Sporadic raids across the frontier marked growing Soviet impatience; and in November, when these pin-pricks failed to impress the Chinese authorities, the Red Army launched an extensive incursion into Chinese territory, dispersing local Chinese forces and capturing two small towns. This time the warning was heeded, and negotiations began in earnest. On December 22, a protocol was signed reinstating the Soviet general manager of the CER and other Soviet officials, restoring the status quo ante, and reserving disputed questions for a future conference. The Red Army had shown up the impotence of the Chinese warlords. The USSR had emerged as a military and diplomatic force in the Far East, and had forged links of common concern with the western Powers. It was a turning-point in Soviet external relations. 
% 注：去除了 "^_ uuai" 的乱码，并将 "status quo ante^" 修正为 "status quo ante,"。

The dispersed, persecuted and dispirited Chinese Communist Party played no part in these events. Instructed by Comintern, its central committee put out the slogan ``Defence of the Soviet Union'', and stepped up its denunciations of the Nanking government. What, however, was seen in Moscow as a flagrant threat to Soviet security appeared to some patriotic Chinese as a move to liberate Chinese territory from foreign, i.e. Soviet, control. Ch'en Tu-hsiu, who had been deposed from the party leadership after the disasters of 1927, was now expelled from the party for voicing these embarrassments, and subsequently proclaimed himself a follower of Trotsky. Here as elsewhere, Comintern could impose discipline, but could not breathe life into the enfeebled party ranks, whose impotence in the urban centres could no longer be disguised. Only Mao Tse-tung's peasant levies and the local Soviets sponsored by them could boast of some successful revolutionary activity. But these exploits were confined to a remote corner of China, and their leaders paid at best hypocritical lip-service to the rulings of the party and of Comintern. Much as the Chinese communist movement owed to Russian inspiration and example, it survived, and would ultimately triumph, in forms unplanned and mistrusted by Moscow.
% 注：修正 "19275" 的 OCR 错误为 "1927,"。